\documentclass[conference]{IEEEtran}
\IEEEoverridecommandlockouts

\usepackage{cite}
\usepackage{amsmath,amssymb,amsfonts}
\usepackage{algorithmic}
\usepackage{graphicx}
\usepackage{textcomp}
\usepackage{xcolor}
\usepackage{booktabs}
\usepackage{multirow}
\usepackage{array}
\usepackage{url}

\graphicspath{{figures/}}

% Space-saving layout (no text changes)
\setlength{\abovecaptionskip}{2pt}
\setlength{\belowcaptionskip}{-6pt}
\setlength{\floatsep}{2pt plus 1pt minus 1pt}
\setlength{\textfloatsep}{4pt plus 2pt minus 2pt}
\setlength{\dblfloatsep}{2pt plus 1pt minus 1pt}
\setlength{\dbltextfloatsep}{4pt plus 2pt minus 2pt}
\setlength{\intextsep}{2pt plus 1pt minus 1pt}
\setlength{\abovedisplayskip}{3pt plus 1pt minus 1pt}
\setlength{\belowdisplayskip}{3pt plus 1pt minus 1pt}
\setlength{\abovedisplayshortskip}{1pt}
\setlength{\belowdisplayshortskip}{3pt plus 1pt minus 1pt}

\def\BibTeX{{\rm B\kern-.05em{\sc i\kern-.025em b}\kern-.08em
    T\kern-.1667em\lower.7ex\hbox{E}\kern-.125emX}}

\begin{document}

\title{Adaptive Confidence-Guided Multi-Scale Vision Transformer
Framework for Explainable Microscopic Fungal Species Classification}

\author{
  \IEEEauthorblockN{Rishi Jayanath A}
  \IEEEauthorblockA{
    \textit{Dept. of Computer Science and Engineering}\\
    \textit{Karunya Institute of Technology and Sciences}\\
    Coimbatore, India\\
    rishijayanath@karunya.edu.in}
  \and
  \IEEEauthorblockN{G. Naveen Sundar}
  \IEEEauthorblockA{
    \textit{Dept. of Computer Science and Engineering}\\
    \textit{Karunya Institute of Technology and Sciences}\\
    Coimbatore, India\\
    naveensundar@karunya.edu}
}

\maketitle

\begin{abstract}
Pathogenic fungi are difficult to identify in the field. The task of reading
microscopic slides is still primarily one that requires training from human
eyes --- a point of delay in clinical decisions and an introduction of
inter-observer error. This paper proposes a new confidence-guided multi-scale
fusion framework called ACMF. The framework runs two Vision Transformer
(ViT-B/16) encoders at different input resolutions --- 224 and 384 pixels ---
and merges the outputs in four different ways. These range from simple
averaging (S1) and accuracy-weighted blending (S2) to a per-image
entropy-based weighting scheme (S3) and an expert MLP fusion head (S4).
Experiments on the DeFungi benchmark --- five clinically relevant fungal
species, 6,801 microscopic images --- demonstrate competitive accuracy.
The best-performing variant is S2, with a macro F1 of 91.63\% on the
held-out test set --- a 3.16-point gain over ViT-B/16-224 used alone.
Attention Rollout maps are generated at inference time, highlighting hyphae
and spore structures without additional annotation. Comparisons with
ResNet-50, EfficientNet-B3, and Swin-Tiny confirm that ACMF delivers the
most consistent value.
\end{abstract}

\begin{IEEEkeywords}
fungal classification, vision transformer, deep learning,
explainable AI, multi-scale fusion, attention rollout, DeFungi, microscopy
\end{IEEEkeywords}

\section{Introduction}

About 1.5 million people per year are killed by fungal infections and have
an impact on more than a billion \cite{kainz2020}. What is particularly
frustrating is that the majority of the species responsible are well-known.
Five of the most prevalent clinical dermatophytes are: \textit{Candida
albicans}, \textit{Aspergillus niger}, \textit{Trichophyton rubrum},
\textit{Trichophyton mentagrophytes}, and \textit{Epidermophyton floccosum}.
Each requires a different antifungal agent, and incorrect diagnosis is not
merely an academic issue --- direct treatment implications follow.

The traditional approach is direct microscopic examination, in which a
trained microscopist examines a stained specimen and distinguishes
morphological clues --- branching pattern of hyphae, the size and arrangement
of spores, the shape of the conidial heads --- to differentiate between
species. This works but is slow, requires years of practice, and experienced
mycologists still disagree on challenging specimens \cite{verweij2015}. Deep
learning has made advances in related problems in pathology
\cite{srinidhi2021}. The subtlety of microscopic fungi has been easy to
overlook: the features to be considered are distributed over a number of
spatial scales at once. Relatively coarse features --- budding yeast clusters
--- are easily visible at normal resolution; fine hyphal wall thicknesses and
septum spacings are only resolvable at higher resolution. A model trained at
a single fixed scale must compromise between these two regimes.

Here, Vision Transformers \cite{dosovitskiy2021} are a logical choice. Global
self-attention captures long-range spatial regularities that vanish from the
vocabulary of convolutional filters, and patch-based tokenisation makes
resolution a tunable parameter. Although even ViT, if used at a single scale,
leaves complementary information on the table --- this is the starting point
for this work.

The Adaptive Confidence-Guided Multi-Scale Fusion (ACMF) framework runs
ViT-B/16 at 224 and 384 pixels, then merges the two sets of class
probabilities. Four strategies are studied --- simple averaging (S1), accuracy
weighting (S2), per-image entropy (S3), and a learned MLP gate (S4) --- so
the comparison is informative about what kind of fusion the task requires.
In addition to accuracy, S3 generates an explicit image-level confidence
signal indicating which resolution was more informative to the model for each
specimen. Attention Rollout maps \cite{abnar2020} ground all predictions in
recognisable morphological structures.

\vspace{-2pt}
\section{Related Work}

\subsection{Bacteria, viruses, fungi and dermatological image classification}

Before deep learning took hold, automated fungal identification relied on
hand-crafted texture descriptors. LBP and GLCM were combined with SVM
classifiers. These pipelines required extensive feature engineering and
generalised poorly across staining protocols. Transfer-learned CNNs changed
this significantly: ResNet \cite{he2016} and EfficientNet \cite{tan2019}
showed that large-scale pretraining followed by task-specific fine-tuning
could be highly effective on skin lesion benchmarks and then on microscopic
fungal datasets. Cassagne et al.'s DeFungi \cite{defungi2021} showed that
transfer-learned CNNs could distinguish the five target species with a fair
degree of accuracy, but confusion between \textit{Aspergillus} and
\textit{Trichophyton} genera continued to be a nuisance --- the gap addressed
directly here.

\subsection{Vision Transformers in Medical Imaging}

In 2020, ViT \cite{dosovitskiy2021} came to be and gave rise to a quick surge
of medical image analysis applications. TransMed \cite{dai2021} adapted
multi-sequence MRI classification using cross-modal attention. Another
architectural strategy employed is the Swin Transformer \cite{liu2021} ---
hierarchical shifted windows that create a spatial pyramid within one model,
handling tasks that require sensitivity to both local texture and global
structure simultaneously. The dual-resolution ViT encoders applied to
microscopic fungal images, and examination of the combined predictions, has
not been reported before in this context.

\subsection{Multi-Scale Fusion and Confidence Weighting}

There is a vast amount of experience with multi-scale strategies in object
detection. Most of the fusion work has been done through feature pyramid
networks \cite{lin2017}, but not in classification at the prediction level
--- that work operates at the feature level. Probability-level ensembling is
less complicated and does not need any sharing between encoders, and ---
importantly --- admits non-parametric confidence weighting; weights are derived
from the model's own prediction distribution, as in multiple-instance learning
\cite{ilse2018}. This idea is extended to the multi-resolution setting with
the S3 strategy.

\subsection{Clinical Models Explainability}

Since then, Grad-CAM \cite{selvaraju2020} has been the most popular
attribution method for medical imaging CNNs. The situation for transformers
is less settled. A principled alternative --- Attention Rollout
\cite{abnar2020} --- propagates multi-head attention across all layers using
residual correction, producing spatially coherent attribution maps without
backpropagation. This gradient-free property makes it practical for
inference-only deployments, which is frequently the case with clinical
software.

\vspace{-2pt}
\section{Dataset and Preprocessing}

\subsection{DeFungi Benchmark}

DeFungi \cite{defungi2021} is a publicly available collection of bright-field
microscopy images from direct mycological examinations, covering five species.
The official split is used throughout: 5,000 training, 899 validation, and
902 test images. Table~\ref{tab:dataset} shows the class distribution. Two
things are worth noting: \textit{C. albicans} (H1) accounts for 40\% of the
training samples while \textit{E. floccosum} (H6) only accounts for 10\%,
resulting in a moderate imbalance that must be addressed by augmentation. For
a few classes (H3, H5, and H6), the total number of specimens in the
validation and test splits is smaller --- a few specimens that still exist in
the original collection.

\begin{table}[!t]
  \caption{DeFungi Dataset Class Distribution}
  \label{tab:dataset}
  \centering
  \footnotesize\setlength{\tabcolsep}{4pt}
  \begin{tabular}{llrrr}
    \toprule
    \textbf{ID} & \textbf{Species} & \textbf{Train} & \textbf{Val} & \textbf{Test}\\
    \midrule
    H1 & \textit{C. albicans}        & 2{,}000 & 437 & 437 \\
    H2 & \textit{A. niger}           &   750   & 233 & 233 \\
    H3 & \textit{T. rubrum}          &   750   &  82 &  82 \\
    H5 & \textit{T. mentagrophytes}  & 1{,}000 &  80 &  80 \\
    H6 & \textit{E. floccosum}       &   500   &  67 &  70 \\
    \midrule
    \multicolumn{2}{l}{\textbf{Total}} & \textbf{5{,}000} & \textbf{899} & \textbf{902}\\
    \bottomrule
  \end{tabular}
\end{table}

\subsection{Augmentation Pipeline}

Training images undergo random resized cropping and horizontal and vertical
flipping (scale 0.8--1.0), with brightness, contrast, and saturation each
$\pm$0.4 (colour jitter). Random erasing \cite{zhong2020} is applied with a
25\% probability, important for occlusion robustness in specimens where debris
blocks part of the field of view. At batch level, MixUp \cite{zhang2018}
($\alpha$ = 0.8) and CutMix \cite{yun2019} ($\alpha$ = 1.0) are applied with
equal probability; both strategies soften inter-class boundaries and have
helped in other fine-grained classification settings with limited data.
Validation and test images see only centre-crop normalisation with ImageNet
statistics. All inputs are independently resized for each encoder at its own
target resolution (224 or 384 pixels) using anti-aliased bilinear
interpolation.

\vspace{-2pt}
\section{Methodology}

\subsection{Backbone Encoders}

Both encoders are initialised from ViT-B/16 \cite{dosovitskiy2021} weights
pretrained on ImageNet-21k, loaded via timm. The original classification head
is replaced with a randomly initialised linear layer mapping the [CLS] token
to 5 class logits. All 12 transformer blocks are fine-tuned end-to-end; no
layers are frozen.

The two encoders share the same architecture but differ in input resolution.
ViT-224 splits each image into $14 \times 14 = 196$ non-overlapping
$16 \times 16$ patches. ViT-384 produces $24 \times 24 = 576$ patches,
providing almost three times the spatial tokens and hence a finer resolution
of the image. Positional embeddings from the ViT-B/16 checkpoint are shared
between the 224- and 384-pixel variants, avoiding the training instability
that can occur when positional encodings are randomly reinitialised at a new
resolution.

\subsection{ACMF Fusion Framework}

Let $\mathbf{p}_{224}$ and $\mathbf{p}_{384}$ be the softmax output vectors
from ViT-224 and ViT-384, both in $\Delta^{C-1}$ for a given image, with
$C = 5$. The four strategies assessed are detailed below.

\textbf{S1 --- Equal-Weight Averaging:}
\begin{equation}
  \hat{\mathbf{p}}^{(1)} = \tfrac{1}{2}\mathbf{p}_{224} + \tfrac{1}{2}\mathbf{p}_{384}
\end{equation}
No assumptions about which encoder is better; the simplest possible baseline
for fusion.

\textbf{S2 --- Accuracy-Weighted Fusion:}
\begin{equation}
  w_{r} = \frac{a_{r}}{a_{224} + a_{384}}, \quad r \in \{224, 384\}
\end{equation}
where $a_r$ is the validation-set top-1 accuracy of encoder $r$. The fused
output becomes $\hat{\mathbf{p}}^{(2)} = w_{224}\,\mathbf{p}_{224} +
w_{384}\,\mathbf{p}_{384}$. This places more trust in the encoder that
generalised better, but the weights are fixed across all test images.

\textbf{S3 --- Confidence-Guided Entropy Weighting:} Rather than assigning
the same weight to every image, S3 computes per-sample weights based on how
confident each encoder is. Confidence is measured through Shannon entropy
$H(\mathbf{p}) = -\sum_{c} p_c \log p_c$: lower entropy means the model is
more certain about its prediction. For each image $i$:
\begin{equation}
  w_{r}^{(i)} = \frac{\bigl(H(\mathbf{p}_{r}^{(i)}) + \epsilon\bigr)^{-1}}
                     {\sum_{r'}\bigl(H(\mathbf{p}_{r'}^{(i)}) + \epsilon\bigr)^{-1}},
                \quad \epsilon = 10^{-8}
\end{equation}
The fused prediction is:
\begin{equation}
  \hat{\mathbf{p}}^{(3)}_i = w_{224}^{(i)}\,\mathbf{p}_{224}^{(i)}
                              + w_{384}^{(i)}\,\mathbf{p}_{384}^{(i)}
\end{equation}
The pair $(w_{224}^{(i)}, w_{384}^{(i)})$ is an interpretable artifact of
inference: it tells which resolution the model found more decisive for
specimen $i$, at no extra computational cost beyond computing entropy of an
already-available softmax vector.

\textbf{S4 --- Learnable MLP Fusion Head:} A two-layer MLP with LayerNorm
and ReLU takes the concatenated vector $[\mathbf{p}_{224};\,\mathbf{p}_{384}]
\in \mathbb{R}^{10}$ and outputs two softmax-normalised mixture weights. The
head is trained for 20 epochs on frozen backbone outputs using AdamW with
cosine annealing, so that the weighting function can in principle learn
arbitrary, specimen-conditioned asymmetries.

\subsection{Training Protocol}

AdamW is used for training, starting at a learning rate of $10^{-4}$, with
weight decay set to $10^{-2}$. The learning rate follows a cosine schedule
after a 5-epoch linear warm-up phase. Label smoothing \cite{szegedy2016} is
applied at $\varepsilon = 0.1$. Automatic mixed precision
(torch.amp.GradScaler) is used throughout training, easing pressure on the
RTX 4050's limited 6~GB VRAM while maintaining numerical stability. The top
checkpoint is preserved by F1 macro, with early stopping triggered after ten
consecutive epochs without a validation F1 gain. Thirty epochs is the upper
limit; ViT-224 converged at epoch~28, and ViT-384 ran all thirty.

\subsection{Attention Rollout for Explainability}

An Attention Rollout map \cite{abnar2020} is generated after every ViT forward
pass. The method multiplies the raw attention through each of the twelve layers
one after another, using a matrix residual correction at each step:
$\hat{A}_l = 0.5\,A_l + 0.5\,I$. The product $\prod_l \hat{A}_l$ gives
per-patch attribution scores relative to the [CLS] token, which are then
restructured into the layout of the patch grid at the original input
resolution. A heat map is overlaid on the original slide image, giving
reviewers a direct visual link between model predictions and the microscopic
structures visible under magnification.

\vspace{-2pt}
\section{Tests With Findings}

All experiments run on a single NVIDIA RTX 4050 Laptop GPU with 6~GB of VRAM,
CUDA 12.9, and PyTorch 2.1. Top-1 accuracy appears alongside macro-averaged
F1 (F1-M), weighted F1, macro-averaged precision, and individual class-level
F1 scores. Among 902 test images, inference latency reflects average GPU
processing time across 1,000 forward passes after 50 warm-up steps. ResNet-50,
EfficientNet-B3, and Swin-Tiny are fine-tuned on DeFungi using the identical
setup applied to the ViT encoders and serve as reference baselines.

\subsection{Performance of Single Models}

Results appear in Table~\ref{tab:baselines}. Leading performance emerges from
Swin-Tiny, with an accuracy of 91.02\% and an F1 macro score of 92.66\%. The
four-level structure manages tasks well due to clear progression through
stages, capturing spatial features at multiple scales within a single forward
pass. ViT-384 trails slightly behind at 91.88\% F1 macro --- the added detail
reveals structural form more fully than details lost when scaled down to 224
pixels.

\begin{table}[!t]
  \caption{Individual Model Performance on DeFungi Test Set}
  \label{tab:baselines}
  \centering
  \footnotesize\setlength{\tabcolsep}{4pt}
  \begin{tabular}{lcccc}
    \toprule
    \textbf{Model} & \textbf{Acc.(\%)} & \textbf{F1-M(\%)} & \textbf{Lat.(ms)} & \textbf{Params(M)}\\
    \midrule
    ResNet-50         & 78.71 & 81.61 &   5.8 & 23.5 \\
    EfficientNet-B3   & 84.26 & 88.04 &  12.7 & 10.7 \\
    ViT-B/16-224      & 84.70 & 88.47 &   8.5 & 85.8 \\
    Swin-Tiny         & \textbf{91.02} & \textbf{92.66} & 11.8 & 27.5 \\
    ViT-B/16-384      & 89.69 & 91.88 & 101.5 & 86.1 \\
    \bottomrule
  \end{tabular}
\end{table}

Close behind, ViT-224 performs about the same as EfficientNet-B3, with results
sitting nearly side by side. ResNet-50 scores lower at 81.61\% F1 --- the gap
stems from missing long-range connections; convolutions process step by step
and miss simultaneous patterns across conidial structures and branching
complexity. A notable trade-off is latency: ViT-384 processes each image in
101.5~ms --- twelve times slower than ViT-224's 8.5~ms --- limiting throughput
to roughly ten images per second. This constraint makes ViT-384 unsuitable for
high-volume screening but acceptable for batch workflows.

\subsection{ACMF Fusion Strategies}

Results from Table~\ref{tab:fusion} alongside Figure~\ref{fig:fusion_bar}
reveal patterns across the four fusion methods. S2 --- adjusted for accuracy
--- shows the highest F1 macro score at 91.63\%, exceeding ViT-224 by 3.16
points yet only 0.25 points ahead of ViT-384 when combining both encoders.
S1 and S3 follow closely at 91.59\% and 91.35\% respectively, both within a
narrow margin. When the MLP is trained, performance drops to 88.59\%.

\begin{table}[!t]
  \caption{ACMF Fusion Strategy Comparison (DeFungi Test Set)}
  \label{tab:fusion}
  \centering
  \footnotesize\setlength{\tabcolsep}{4pt}
  \begin{tabular}{lcccc}
    \toprule
    \textbf{Strategy} & \textbf{Acc.(\%)} & \textbf{F1-M(\%)} & \textbf{F1-W(\%)} & \textbf{Prec-M(\%)}\\
    \midrule
    S1 Equal Weight      & 89.47 & 91.59 & 89.50 & 90.86 \\
    S2 Accuracy Weighted & 89.36 & \textbf{91.63} & 89.39 & 90.80 \\
    S3 Confidence Guided & 89.25 & 91.35 & 89.29 & 90.71 \\
    S4 Learnable MLP     & 84.92 & 88.59 & 85.07 & 87.96 \\
    \midrule
    ViT-224 (ref.)       & 84.70 & 88.47 & 84.86 & 87.84 \\
    ViT-384 (ref.)       & 89.69 & 91.88 & 89.70 & 91.41 \\
    \bottomrule
  \end{tabular}
\end{table}

\begin{figure}[!t]
  \vspace{-4pt}
  \centering
  \includegraphics[width=0.92\columnwidth]{fusion_comparison.png}
  \caption{F1 macro and accuracy scores for all four ACMF fusion strategies,
           each shown in a distinct colour.}
  \label{fig:fusion_bar}
  \vspace{-4pt}
\end{figure}

The narrow gap between S1, S2, and S3 --- just 0.28 points --- is instructive
on its own. The primary benefit of the ACMF framework comes from combining two
resolutions entirely, not from the specific weighting mechanism. The three
heuristic approaches rely mostly on practical rules and perform similarly on
macro F1, a sign that improvements hold steady even when conditions shift.
That consistency suggests strength in design, not fine-tuned luck in adjusting
the weight assignment method.

S4's weaker result calls for a closer look. Though most expressive by nature,
the MLP head is limited to just 5,000 training samples when learning to combine
softmax vectors from a dual-encoder setup. The input space is ten-dimensional
--- two softmax vectors over five classes --- and the calibration patterns
between encoders can vary. This makes optimising the MLP head challenging.
Despite limited information, S3 uses entropy as a guiding principle entirely
without instruction, avoiding assumptions about data patterns, yet still matches
the performance of S1 and S2, producing a clear confidence measure for each
image as a byproduct.

Examining the recorded weight pairs $(w^{(i)}_{224}, w^{(i)}_{384})$ across
the evaluation data is instructive. ViT-224 received consistently greater
importance --- less uncertainty --- for H5 and H6 samples, which show traits
best seen at larger scales appearing clearly under regular magnification.
ViT-384 dominated for H2 and H3, where delicate hyphal wall structures and
septum spacings emerge only at fine detail. That the model's behaviour matches
what an expert microscopist would look for is a form of quiet validation.

\subsection{Per-Class Analysis}

Looking at Table~\ref{tab:perclass}, outcomes are sorted by class, revealing
a consistent trend: H2 (\textit{A. niger}) consistently ranks as the most
difficult category, while H5 and H6 stand out as simplest. ResNet-50's H2 F1
score stays notably low at 60.09\% --- models built on newer architectures
perform significantly stronger, with ViT-384 hitting 82.80\% and Swin-Tiny
reaching 85.22\%. Global context appears important for identifying
\textit{A. niger}. The challenge is notable: \textit{A. niger} specimens in
this dataset show substantial variability in conidial head dimensions, with
branching patterns and angles shifting across specimens, making these features
more elusive to identify compared to the four other classes.

\begin{table}[!t]
  \caption{Per-Class F1 Score (\%) --- DeFungi Test Set}
  \label{tab:perclass}
  \centering
  \footnotesize\setlength{\tabcolsep}{3pt}
  \begin{tabular}{lcccccc}
    \toprule
    \textbf{Model} & \textbf{H1} & \textbf{H2} & \textbf{H3} & \textbf{H5} & \textbf{H6} & \textbf{F1-M}\\
    \midrule
    ResNet-50        & 82.96 & 60.09 & 79.78 & 92.59 & 92.65 & 81.61 \\
    EfficientNet-B3  & 85.58 & 73.49 & 88.89 & 94.41 & 97.81 & 88.04 \\
    ViT-B/16-224     & 85.03 & 75.60 & 91.12 & 95.65 & 94.96 & 88.47 \\
    ViT-B/16-384     & 90.40 & 82.80 & 92.77 & 96.25 & 97.18 & 91.88 \\
    Swin-Tiny        & 91.82 & 85.22 & 93.75 & 96.86 & 95.65 & 92.66 \\
    \midrule
    S2 ACMF          & 90.06 & 82.43 & 93.41 & 95.06 & 97.18 & \textbf{91.63}\\
    \bottomrule
  \end{tabular}
\end{table}

Notably, S2 fusion achieves 93.41\% on H3 (\textit{T. rubrum}), higher than
each separate encoder at 92.77\% and 91.12\% respectively --- a clear case of
merging elements offering real improvement compared to each part alone.
Regarding H2, S2 achieves 82.43\%, a touch lower than ViT-384 alone, yet
still outperforming ViT-224 by a notable margin.

\subsection{Explainability Results}

Attention rollout maps on the test set show where correct predictions link to
where attention settles in space. For H1 samples, focus shifts to sequences
of yeast cells forming buds; for H2 and H3, attention settles on hyphal
surfaces and at junctions where strands split; for H5 and H6, focus shifts
toward the pointed ends of macroconidia and club-shaped terminal cells. For
H6 correct outcomes, rollout mapping typically gives negligible weight to
background hyphae, shifting attention almost entirely to the tip where it
thickens slightly --- the precise marker a mycologist would seek during
examination.

Wrong labels reveal something else. Maps for wrongly-labelled H2 samples
appear nearly uniform across the image, with attention blended throughout
without distinct focus. These samples may not present the clear structural
features the model learned to rely on --- a situation a skilled microscopist
would also find challenging. Cases that show strong softmax signals alongside
scattered spatial attention patterns may warrant flagging for expert review,
even when the model appears confident, given the prediction fails to rely on
the typical distinguishing characteristics.

\vspace{-2pt}
\section{Discussion}

The result that stands out most is Swin-Tiny outperforming all four ACMF
strategies despite using a single model and a single forward pass. Performance
reaches 92.66\% on F1 macro, slightly above 91.63\% --- a small but consistent
edge across every measure. The explanation is architectural: Swin's layered
window attention creates a spatial pyramid internally, capturing features at
multiple sizes simultaneously. ACMF handles this through execution --- one
encoder operates at a set resolution, the other does likewise, before the
outputs merge at the probability level. Internal representations carry more
information than what emerges through post-hoc likelihood weighting.

Still, this should not be interpreted as proof that ACMF lacks practical
value. Working across multiple resolutions consistently helps performance
compared to a basic ViT setup with a fixed scale. In settings where ViT is
already deployed --- given established infrastructure, access to pretrained
models, and the interpretability that Attention Rollout brings --- combining
dual resolutions through fusion is an approach that works within existing
structures without demanding redesign.

The comparison among S1, S2, S3, and S4 tells a useful story about the kind
of blending this task requires. Similar performance across the three heuristic
approaches indicates that handling the task does not require detailed per-image
adjustments on this scale of data. When S4 falls short, it echoes a familiar
lesson: learned ensemble weighting is more data-hungry than heuristic methods.
Five thousand training samples fall short for the MLP to learn calibration
patterns beyond what the frozen encoders already provide.

\vspace{-2pt}
\section{Conclusion}

The ACMF framework is presented with an assessment of four strategies for
combining ViT-B/16 predictions at two input resolutions on the DeFungi
microscopic fungal classification benchmark. The accuracy-weighted strategy
(S2) reaches an F1 macro of 91.63\%, rising by 3.16 points compared to
ViT-224 alone and sitting just a quarter-point beneath ViT-384, which runs
twelve times slower. The entropy-guided method (S3) achieves similar
performance while generating per-image confidence indicators that signal to
an operator which resolution the model found more informative. Rollout patterns
across the evaluation set confirm that model predictions tie back to the
morphological structures --- hyphal form, spore shape, and branching geometry
--- that count toward trust in clinical settings.

Constraints are worth noting. Running dual encoders cuts inference speed by
roughly half compared to a single model, making offline or batch settings the
main practical scenarios. Performance dips slightly below Swin-Tiny despite
the lower parameter count of the heuristic fusion methods. The trained S4 head
faces restrictions due to limited data --- 5,000 training samples fall short
for learning calibration patterns beyond what the heuristic approaches already
provide.

Future possibilities appear most compelling when combining a Swin-based
encoder with a ViT encoder rather than two ViTs --- varied building styles
could improve resilience and capture complementary representations. Distilling
insight from the dual-encoder pipeline into a lightweight single framework
might uncover performance gains without added latency. Data size still holds
things back even as models improve --- pretraining on unlabelled clinical
microscopy data could expand the useful pool of learning examples without
requiring additional expert annotation.

\section*{Acknowledgment}

The Department of Computer Science and Engineering at Karunya Institute of
Technology and Sciences provided the computing facilities and academic
environment for this work. Acknowledgment is also due to the developers of
PyTorch, timm, and the DeFungi dataset for making these resources publicly
available. This work received no external funding.

\vspace{-2pt}
{\small
\begin{thebibliography}{00}

\bibitem{kainz2020}
K. Kainz, M. Lakemeyer, and F. Madeo,
``Fungal diseases --- forgotten killers? Understanding the global burden
of fungal infections,''
\textit{EMBO Mol. Med.}, vol. 12, no. 8, p. e12114, 2020.

\bibitem{verweij2015}
P. E. Verweij and J. F. Meis,
``Microbiological diagnosis of invasive fungal infections: challenges and
progress,''
\textit{J. Antimicrob. Chemother.}, vol. 55, suppl. 1, pp. 15--22, 2005.

\bibitem{srinidhi2021}
C. L. Srinidhi, O. Ciga, and A. L. Martel,
``Deep neural network models for computational histopathology: a survey,''
\textit{Med. Image Anal.}, vol. 67, p. 101813, 2021.

\bibitem{dosovitskiy2021}
A. Dosovitskiy \textit{et al.},
``An image is worth 16x16 words: Transformers for image recognition at scale,''
in \textit{Proc. ICLR}, 2021.

\bibitem{he2016}
K. He, X. Zhang, S. Ren, and J. Sun,
``Deep residual learning for image recognition,''
in \textit{Proc. IEEE CVPR}, pp. 770--778, 2016.

\bibitem{tan2019}
M. Tan and Q. V. Le,
``EfficientNet: Rethinking model scaling for convolutional neural networks,''
in \textit{Proc. ICML}, vol. 97, pp. 6105--6114, 2019.

\bibitem{liu2021}
Z. Liu \textit{et al.},
``Swin transformer: Hierarchical vision transformer using shifted windows,''
in \textit{Proc. IEEE ICCV}, pp. 10012--10022, 2021.

\bibitem{abnar2020}
S. Abnar and W. Zuidema,
``Quantifying attention flow in transformers,''
in \textit{Proc. ACL}, pp. 4190--4197, 2020.

\bibitem{selvaraju2020}
R. R. Selvaraju \textit{et al.},
``Grad-CAM: Visual explanations from deep networks via gradient-based
localization,''
\textit{Int. J. Comput. Vis.}, vol. 128, no. 2, pp. 336--359, 2020.

\bibitem{defungi2021}
P. Cassagne \textit{et al.},
``DeFungi: Direct mycological examination of microscopic fungi images,''
\textit{arXiv preprint arXiv:2110.04984}, 2021.

\bibitem{lin2017}
T.-Y. Lin, P. Doll\'{a}r, R. Girshick, K. He, B. Hariharan, and S. Belongie,
``Feature pyramid networks for object detection,''
in \textit{Proc. IEEE CVPR}, pp. 2117--2125, 2017.

\bibitem{ilse2018}
M. Ilse, J. M. Tomczak, and M. Welling,
``Attention-based deep multiple instance learning,''
in \textit{Proc. ICML}, vol. 80, pp. 2127--2136, 2018.

\bibitem{zhang2018}
H. Zhang, M. Cisse, Y. N. Dauphin, and D. Lopez-Paz,
``MixUp: Beyond empirical risk minimization,''
in \textit{Proc. ICLR}, 2018.

\bibitem{yun2019}
S. Yun, D. Han, S. J. Oh, S. Chun, J. Choe, and Y. Yoo,
``CutMix: Training strategy that makes strongly localized contributions matter,''
in \textit{Proc. IEEE ICCV}, pp. 6023--6032, 2019.

\bibitem{zhong2020}
Z. Zhong, L. Zheng, G. Kang, S. Li, and Y. Yang,
``Random erasing data augmentation,''
in \textit{Proc. AAAI}, vol. 34, pp. 13001--13008, 2020.

\bibitem{szegedy2016}
C. Szegedy, V. Vanhoucke, S. Ioffe, J. Shlens, and Z. Wojna,
``Rethinking the Inception architecture for computer vision,''
in \textit{Proc. IEEE CVPR}, pp. 2818--2826, 2016.

\bibitem{dai2021}
Y. Dai \textit{et al.},
``TransMed: Transformers advance multi-modal medical image classification,''
\textit{Diagnostics}, vol. 11, no. 8, p. 1384, 2021.

\end{thebibliography}
} % end \small

\end{document}
